% Blueprint: compactness-outline
% Created by Numina

\begin{document}

\section{Worms and thickenings}

\begin{definition}[Worm]
    \label{def:worm}
    A \emph{worm} is a $1$-Lipschitz curve $f : [0,1] \to \RR^2$. A worm is
    \emph{pinned} if $f(0) = (0,0)$. Since $f$ is $1$-Lipschitz on $[0,1]$,
    a pinned worm satisfies $\lvert f(x) \rvert \le x \le 1$ for all
    $x \in [0,1]$.
\end{definition}

\begin{definition}[$\varepsilon$-thickening]
    \label{def:thickening}
    For a set $A \subseteq \RR^2$ and $\varepsilon \ge 0$, the
    \emph{$\varepsilon$-thickening} of $A$ is the closed
    $\varepsilon$-neighbourhood
    \[
        A^{\varepsilon} \;=\; \{\, x \in \RR^2 : \dist(x, A) \le \varepsilon \,\}.
    \]
    We say a worm $f$ is \emph{$\varepsilon$-close} to a worm $g$ if
    $\sup_{x \in [0,1]} \lvert f(x) - g(x) \rvert \le \varepsilon$, which
    implies that the image of $f$ is contained in the
    $\varepsilon$-thickening of the image of $g$.
\end{definition}

\begin{definition}[Grid polygonal worm]
    \label{def:gridPolygonalWorm}
    Fix a step $\delta > 0$ and a vertex count $k \in \NN$. A
    \emph{$(k,\delta)$-grid polygonal worm} is a pinned worm obtained by
    choosing grid points $p_0, p_1, \dots, p_k \in \delta\ZZ \times \delta\ZZ$
    with $p_0 = (0,0)$, setting $\tilde f(n/k) = p_n$ for $0 \le n \le k$, and
    interpolating linearly on each interval $[n/k, (n+1)/k]$, subject to the
    resulting curve being $1$-Lipschitz.
\end{definition}

\section{A finite \texorpdfstring{$\varepsilon$}{epsilon}-net of polygonal worms}

\begin{theorem}[Rounding a worm to the grid]
    \label{thm:roundedWorm}
    \uses{def:worm, def:gridPolygonalWorm, def:thickening}
    Let $f : [0,1] \to \RR^2$ be a pinned $1$-Lipschitz worm, and let
    $k \in \NN$ and $\delta > 0$. Define
    \[
        f_1 = \frac{1}{1 + 2\delta k}\, f,
    \]
    let $\tilde f(n/k)$ be a nearest point of $\delta\ZZ \times \delta\ZZ$ to
    $f_1(n/k)$ for $0 \le n \le k$, and extend $\tilde f$ by linear
    interpolation. Then $\tilde f$ is a $(k,\delta)$-grid polygonal worm and
    \[
        \sup_{x \in [0,1]} \lvert f(x) - \tilde f(x) \rvert
            \;\le\; 2\delta k + \delta + \frac{1}{k}.
    \]
\end{theorem}
\begin{proof}
    \uses{def:worm, def:gridPolygonalWorm}
    The rescaling $f_1 = f / (1 + 2\delta k)$ is
    $\tfrac{1}{1+2\delta k}$-Lipschitz. Because $f$ is pinned,
    $\lvert f(x) \rvert \le 1$, so
    \[
        \lvert f(x) - f_1(x) \rvert
            = \lvert f(x) \rvert \cdot \frac{2\delta k}{1 + 2\delta k}
            \le 2\delta k .
    \]
    Each node $f_1(n/k)$ is replaced by a nearest grid point, moving it by at
    most $\delta$ (half a grid cell in each coordinate). Hence the nodes of
    $\tilde f$ are within $\delta$ of those of $f_1$, and on each subinterval of
    length $1/k$ the segment of $\tilde f$ has length at most
    $\tfrac{1}{k}\cdot\tfrac{1}{1+2\delta k} + 2\delta$; multiplying by $k$
    shows the Lipschitz constant of $\tilde f$ is at most
    $\tfrac{1}{1+2\delta k}\,(1 + 2\delta k) = 1$, so $\tilde f$ is a worm.
    Finally, comparing $f$ to $\tilde f$ at an arbitrary $x \in [0,1]$, the
    rescaling contributes at most $2\delta k$, the grid rounding at most
    $\delta$, and the linear interpolation between nodes at most $1/k$,
    giving the stated bound.
\end{proof}

\begin{theorem}[Finite $\varepsilon$-net of worms]
    \label{thm:finiteEpsilonNet}
    \uses{def:worm, def:gridPolygonalWorm, def:thickening, thm:roundedWorm}
    For every $\varepsilon > 0$ there is a finite set $S_\varepsilon$ of pinned
    worms such that every pinned worm is $\varepsilon$-close to some element of
    $S_\varepsilon$; equivalently, its image is contained in the
    $\varepsilon$-thickening of the image of some element of $S_\varepsilon$.
    Moreover $S_\varepsilon$ may be taken to consist of $(k,\delta)$-grid
    polygonal worms, and is computable.
\end{theorem}
\begin{proof}
    \uses{thm:roundedWorm}
    Choose $k > 2/\varepsilon$, so that $1/k < \varepsilon/2$, and then choose
    $\delta > 0$ small enough that $2\delta k + \delta < \varepsilon/2$. Let
    $S_\varepsilon$ be the set of all $(k,\delta)$-grid polygonal worms whose
    nodes lie within distance $1$ of the origin; this is a finite set, since
    each of the finitely many nodes ranges over the finitely many grid points
    of $\delta\ZZ \times \delta\ZZ$ inside the unit disc. Given any pinned worm
    $f$, Theorem~\ref{thm:roundedWorm} produces a $(k,\delta)$-grid polygonal
    worm $\tilde f \in S_\varepsilon$ with
    \[
        \sup_x \lvert f(x) - \tilde f(x) \rvert
            \le 2\delta k + \delta + \frac{1}{k} < \varepsilon .
    \]
    Hence $f$ is $\varepsilon$-close to $\tilde f$, and its image lies in the
    $\varepsilon$-thickening of the image of $\tilde f$. Finiteness and the
    explicit construction make $S_\varepsilon$ computable.
\end{proof}

\section{Bounds on the Moser number}

\begin{definition}[Moser cover number]
    \label{def:moserNumber}
    \uses{def:worm}
    A set $C \subseteq \RR^2$ is a \emph{cover} if for every worm $f$ there is
    an orientation-preserving isometry (a translation--rotation) $g$ of $\RR^2$
    such that the image of $g \circ f$ is contained in $C$. The \emph{Moser
    cover number} $M$ is the infimum of $\area(C)$ over all convex covers $C$.
\end{definition}

\begin{definition}[Minimal convex cover of a finite set of worms]
    \label{def:minimalCover}
    \uses{def:worm}
    For a finite set $S$ of worms, the \emph{minimal convex cover} $K(S)$ is a
    convex set of least area for which there exist translation--rotations
    $(g_s)_{s \in S}$ with the image of $g_s \circ s$ contained in $K(S)$ for
    every $s \in S$; that is, $K(S)$ minimises
    $\area\!\big(\conv \bigcup_{s \in S} g_s(\image\, s)\big)$ over all choices
    of translation--rotations $g_s$.
\end{definition}

\begin{remark}
    \label{rem:computeMinimalCover}
    \uses{def:minimalCover}
    The minimal convex cover $K(S)$ is the minimum, over all
    translation--rotations of the elements of $S$, of the area of the convex
    hull of their union. For a fixed finite $S$ this is a finite-dimensional
    optimisation problem in the placement parameters and is in principle
    computable, e.g.\ by quantifier elimination over the reals, and plausibly
    by convex optimisation.
\end{remark}

\begin{theorem}[Lower bound on the Moser number]
    \label{thm:lowerBound}
    \uses{def:moserNumber, def:minimalCover}
    For any finite set $S$ of worms, $M \ge \area\big(K(S)\big)$.
\end{theorem}
\begin{proof}
    \uses{def:moserNumber, def:minimalCover}
    Let $C$ be any convex cover. Since every element of $S$ is a worm, for each
    $s \in S$ there is a translation--rotation $g_s$ placing $g_s \circ s$
    inside $C$. Thus $C$ is a convex set admitting placements of all elements of
    $S$, so by minimality $\area(C) \ge \area(K(S))$. Taking the infimum over
    convex covers $C$ gives $M \ge \area(K(S))$.
\end{proof}

\begin{theorem}[Upper bound on the Moser number]
    \label{thm:upperBound}
    \uses{def:moserNumber, def:minimalCover, def:thickening, thm:finiteEpsilonNet}
    Let $\varepsilon > 0$ and let $S_\varepsilon$ be a finite
    $\varepsilon$-net of worms as in Theorem~\ref{thm:finiteEpsilonNet}. Then
    \[
        M \le \area\!\big(K(S_\varepsilon)^{\varepsilon}\big),
    \]
    the area of the $\varepsilon$-thickening of the minimal convex cover of
    $S_\varepsilon$.
\end{theorem}
\begin{proof}
    \uses{def:moserNumber, def:minimalCover, def:thickening, thm:finiteEpsilonNet}
    Write $K = K(S_\varepsilon)$ with placements $(g_s)_{s \in S_\varepsilon}$.
    The $\varepsilon$-thickening $K^{\varepsilon}$ is convex. We claim it is a
    cover. Let $f$ be an arbitrary worm; after a translation--rotation we may
    assume $f$ is pinned. By Theorem~\ref{thm:finiteEpsilonNet} there is
    $s \in S_\varepsilon$ and a translation--rotation $h$ with $f$
    $\varepsilon$-close to $h \circ s$. Then $g_s \circ h^{-1}$ maps the image
    of $f$ into the $\varepsilon$-thickening of the image of $g_s \circ s$,
    which lies in $K^{\varepsilon}$. Hence $K^{\varepsilon}$ is a convex cover,
    and $M \le \area(K^{\varepsilon})$.
\end{proof}

\begin{theorem}[Steiner gap between the bounds]
    \label{thm:steinerGap}
    \uses{def:minimalCover, def:thickening}
    Let $K$ be a convex body in $\RR^2$ with perimeter $L$. Then for every
    $\varepsilon \ge 0$,
    \[
        \area\big(K^{\varepsilon}\big)
            = \area(K) + \varepsilon\, L + \pi \varepsilon^{2}.
    \]
    In particular the upper and lower bounds of
    Theorems~\ref{thm:lowerBound} and~\ref{thm:upperBound} differ by at most
    $\varepsilon\, L + \pi \varepsilon^{2}$, where $L$ is the perimeter of
    $K(S_\varepsilon)$.
\end{theorem}
\begin{proof}
    \uses{def:thickening}
    This is the planar Steiner formula for convex bodies: the outer parallel
    body at distance $\varepsilon$ adds a strip of area $\varepsilon L$ along
    the boundary together with disc sectors of total area $\pi \varepsilon^2$ at
    the convex corners.
\end{proof}

\begin{remark}[Perimeter bound]
    \label{rem:perimeterBound}
    \uses{def:minimalCover, thm:steinerGap}
    To turn Theorem~\ref{thm:steinerGap} into an effective error bound one needs
    an a priori bound on the perimeter $L$ of $K(S_\varepsilon)$. Insisting that
    $S_\varepsilon$ contains a fixed closed worm (a loop, e.g.\ a small circle)
    forces every cover to be bounded, hence $K(S_\varepsilon)$ is bounded with a
    controlled perimeter. A quantitative perimeter bound is left as a
    \textsc{todo}; one may also add a square to $S_\varepsilon$ to make this
    argument go through.
\end{remark}

\begin{theorem}[Approximating the Moser number]
    \label{thm:approxAlgorithm}
    \uses{def:moserNumber, thm:finiteEpsilonNet, thm:lowerBound, thm:upperBound, thm:steinerGap, rem:perimeterBound}
    Assume a perimeter bound $L_0$ on $K(S_\varepsilon)$ as in
    Remark~\ref{rem:perimeterBound}. Then for any target accuracy $x > 0$ the
    Moser cover number $M$ can be computed to within $x$ by the following
    procedure:
    \begin{enumerate}
        \item choose $\varepsilon = x$ small enough that
            $\varepsilon L_0 + \pi \varepsilon^2 \le x$;
        \item compute a finite $\varepsilon$-net $S_\varepsilon$, for instance
            by repeatedly locating worms not yet covered and adding them;
        \item compute the minimal convex cover $K(S_\varepsilon)$.
    \end{enumerate}
    Then $\area(K(S_\varepsilon))$ is a lower bound and
    $\area(K(S_\varepsilon)^{\varepsilon})$ an upper bound for $M$, and the two
    differ by at most $x$.
\end{theorem}
\begin{proof}
    \uses{thm:lowerBound, thm:upperBound, thm:steinerGap}
    By Theorem~\ref{thm:lowerBound}, $\area(K(S_\varepsilon)) \le M$, and by
    Theorem~\ref{thm:upperBound}, $M \le \area(K(S_\varepsilon)^{\varepsilon})$.
    By Theorem~\ref{thm:steinerGap} the gap between these bounds equals
    $\varepsilon L + \pi \varepsilon^2 \le \varepsilon L_0 + \pi \varepsilon^2
    \le x$. Hence either bound approximates $M$ to within $x$.
\end{proof}

\end{document}
